\chapter{ICNN, Brenier Maps, and Monge-Gap Direct Map Learning}
\label{ch:bf-neural-ot-direct-map}

The previous chapter studied the narrowest neural extension of the entropic OT
teacher: retain the teacher and learn how to start its solver better.  The
present chapter changes the question.  Instead of accelerating a declared
solver, it asks how one can learn the transport map itself.

This is a more ambitious route.  The potential gain is obvious: if a good map is
learned directly, one may bypass repeated iterative solves at deployment time.
The cost is equally important: the learned object is no longer an initialization
for a retained teacher.  It is the transport map itself, so every modeling and
training choice now acts directly on the output cloud.

This chapter therefore has one dominant task: identify what object the direct-map
lane is trying to learn, why Brenier/ICNN and Monge-gap methods are mathematically
natural representatives of that lane, and what evidence is needed once no
corrective teacher solve is left in control.

\section{What problem changes in this chapter?}
\label{sec:bf-neural-ot-direct-orientation}

The best way to read this chapter is to compare it with the retained-teacher
warm-start chapter.
\begin{itemize}
  \item In the warm-start chapter, the teacher still decides what the final
  transport plan and output cloud are; the learner only helps the teacher solve
  its own problem faster.
  \item In the present chapter, the learner itself proposes the transport map.
  There may be no corrective Sinkhorn phase at all.
\end{itemize}
So the central question is no longer ``did we accelerate the same teacher?''  It
is ``what map object have we learned, and under what assumptions should we trust
it?''

\paragraph{Plain-language takeaway.}
If the previous chapter was about helping an expensive trusted calculation, this
chapter is about replacing that calculation by a learned map.  That is why the
mathematical structure and the evidence contract both change.

\section{Running direct-map cell}
\label{sec:bf-neural-ot-direct-running-cell}

A tiny one-dimensional cell is enough to carry the main ideas.  Let the source
be the two-point cloud
\[
  X_{\mathrm{src}}=\{-1,1\}
  \quad\text{with equal weights }\frac{1}{2},
\]
and let the target be
\[
  X_{\mathrm{tgt}}=\{0,2\}
  \quad\text{with equal weights }\frac{1}{2}.
\]
The obvious translation map is
\begin{equation}
\label{eq:bf-neural-ot-direct-shift-map}
  T^\star(x)=x+1.
\end{equation}
Applying \(T^\star\) to the source cloud sends \(-1\mapsto 0\) and
\(1\mapsto 2\), so the pushforward equals the target cloud exactly.

For quadratic cost,
\begin{equation}
\label{eq:bf-neural-ot-direct-shift-cost}
  \frac{1}{2}\sum_{x\in X_{\mathrm{src}}}|x-T^\star(x)|^2
  =
  \frac{1}{2}(1^2+1^2)=1.
\end{equation}
This simple cell already shows the difference between a direct map and a retained
teacher solve.  The output cloud is obtained directly by evaluating \(T\) on the
source particles; there is no intermediate coupling matrix or balancing loop in
this picture.

A second map on the same endpoint cloud shows why direct-map evidence is more
subtle.  Consider
\begin{align}
\label{eq:bf-neural-ot-direct-good-bad-maps}
  T_{\mathrm{good}}(-1)&=0,
  &T_{\mathrm{good}}(1)&=2,\\
  T_{\mathrm{bad}}(-1)&=2,
  &T_{\mathrm{bad}}(1)&=0.
\end{align}
Both maps reach the same endpoint set \(\{0,2\}\), but for quadratic cost the
good map has cost
\[
  \frac{1}{2}(1^2+1^2)=1,
\]
whereas the bad map has cost
\[
  \frac{1}{2}(3^2+1^2)=5.
\]
So endpoint fit alone does not settle whether the learned map behaves like a
good transport map for the declared cost.

\section{Direct-map lane: object changed, evidence changed}
\label{sec:bf-neural-ot-direct-object}

A direct map-learning method tries to produce a deterministic transport map
\begin{equation}
\label{eq:bf-neural-ot-direct-map-object}
  T:
  \mathbb R^{d_x}\to\mathbb R^{d_x}.
\end{equation}
The source measure \(\mu\) and target measure \(\nu\) are related by the
pushforward relation
\begin{equation}
\label{eq:bf-neural-ot-direct-pushforward}
  T_\#\mu = \nu,
\end{equation}
meaning that if one samples \(X\sim\mu\) and then computes \(T(X)\), the law of
that transformed sample is \(\nu\).

The important change relative to the warm-start chapter is immediate: the learned
object is no longer an internal solver coordinate.  It is the transport map
itself, or a potential from which that map is derived.  In the vocabulary of the
OT block, retained-teacher learning changes the execution path while a corrected
solve still returns to a declared transport object.  Direct-map learning changes
the transport object directly and therefore inherits a different evidence
contract.

\begin{center}
\small
\begin{tabular}{p{0.24\textwidth} p{0.23\textwidth} p{0.41\textwidth}}
\toprule
Quantity & Status in this chapter & Practical burden \\
\midrule
learned map \(T\) & primary learned object & must be validated directly, not through a retained teacher residual \\
pushforward \(T_\#\mu\) & endpoint-law target & should match the declared target law or cloud closely enough on held-out problems \\
transport-cost / Monge-gap behavior & auxiliary OT-structure evidence & helps distinguish a plausible transport map from an arbitrary endpoint fit \\
downstream scalar built from mapped particles & implementation consequence & must be checked for drift once the map is embedded in filtering or HMC-adjacent code \\
\bottomrule
\end{tabular}
\end{center}

\paragraph{Reader checkpoint.}
The direct-map route is not just a faster warm-start.  It asks the learned map
itself to carry the approximation burden.

\section{Brenier map viewpoint}
\label{sec:bf-neural-ot-direct-brenier}

In the quadratic-cost Euclidean setting, a classical route is to represent the
map as the gradient of a convex potential:
\begin{equation}
\label{eq:bf-neural-ot-direct-brenier}
  T^\star(x)=\nabla\psi^\star(x),
  \qquad
  \psi^\star\text{ convex}.
\end{equation}
This statement should be read with its theorem-level assumptions attached.  The
classical Brenier characterization is a continuous-measure result for quadratic
cost in Euclidean space under regularity conditions such as finite second
moments and absolute continuity of the source measure with respect to Lebesgue
measure.  In the finite empirical-cloud setting used throughout BayesFilter, a
Monge map may fail to exist or fail to be unique in that theorem-strength sense.
So the present chapter uses the Brenier viewpoint primarily as structured
motivation for a learned map family, not as a blanket theorem that every finite
cloud transport problem in the particle-filter setting is exactly represented by
a gradient of a convex potential.

In one dimension, this means the derivative of the potential defines the map.
For the translation cell above, one convenient choice is
\begin{equation}
\label{eq:bf-neural-ot-direct-brenier-example}
  \psi^\star(x)=\frac{1}{2}x^2+x.
\end{equation}
Then
\begin{equation}
\label{eq:bf-neural-ot-direct-brenier-example-grad}
  \nabla\psi^\star(x)=x+1=T^\star(x).
\end{equation}
The second derivative of \(\psi^\star\) is \(1\), so the potential is convex.
This tiny calculation is the cleanest way to see what the Brenier/ICNN family is
trying to exploit: instead of learning the map directly, it learns a structured
potential whose gradient gives the map.

\section{ICNN/Brenier route: structured direct-map implementation}
\label{sec:bf-neural-ot-direct-icnn}

Input-convex neural networks (ICNNs) provide a parametric family of scalar
functions that remain convex in designated inputs.  The implementation idea is:
\begin{enumerate}[label=(\roman*)]
  \item parameterize the potential \(\psi\) by an ICNN \(\psi_\theta\),
  \item obtain the transport map by differentiating the network,
  \(
    T_\theta(x)=\nabla\psi_\theta(x),
  \)
  \item train \(\psi_\theta\) so that the induced map sends source mass toward
  the target according to the chosen OT objective.
\end{enumerate}

A useful objective to keep in view is the convex-potential dual form
\begin{equation}
\label{eq:bf-neural-ot-direct-dual-objective}
  J(\psi)
  =
  \mathbb E_{X\sim\mu}[\psi(X)]
  +
  \mathbb E_{Y\sim\nu}[\psi^*(Y)],
\end{equation}
where \(\psi^*\) is the convex conjugate
\begin{equation}
\label{eq:bf-neural-ot-direct-conjugate}
  \psi^*(y)=\sup_x\{x^\top y-\psi(x)\}.
\end{equation}
This is best read as continuous-measure theorem language or as motivation for a
structured parametric family, not as a claim that an arbitrary finite-cloud
training loop automatically inherits the full optimal-transport theorem.  In a
finite empirical implementation, one still has to choose how \(\mu\) and \(\nu\)
are represented, how gradients of the potential are evaluated on particle clouds,
and how any conjugate or dual approximation is actually computed numerically.
The attraction of this route is that the map is not a black-box regressor.  It is
constrained to come from a convex potential, which is exactly the sort of
structure the quadratic-cost theory suggests under its stated assumptions.

The cost is also immediate from the equations.  One must evaluate gradients of
the potential and, depending on the chosen training formulation, may need to
evaluate or approximate the convex conjugate.  So the route trades a retained
transport solver for a more structured but more delicate map-learning objective.

\begin{algorithm}[ht]
\caption{\texttt{train\_icnn\_brenier\_map}}
\begin{algorithmic}[1]
\Require source-cloud sampler for \(\mu\), target-cloud sampler for \(\nu\), ICNN potential family \(\psi_\theta\), declared quadratic-cost/Brenier setting, convexity-enforcement policy, optimizer, learning-rate schedule, batch size, stopping rule, and held-out validation clouds
\Ensure trained potential \(\psi_\theta\), induced map \(T_\theta(x)=\nabla\psi_\theta(x)\), and validation diagnostics for pushforward mismatch and scalar drift
\State Initialize ICNN parameters \(\theta\) and optimizer state
\While{the stopping rule has not been met}
  \State Draw minibatches \(x_b\sim\mu\) and \(y_b\sim\nu\)
  \State Evaluate the potential values \(\psi_\theta(x_b)\)
  \State Evaluate the gradient-defined transport map \(T_\theta(x_b)=\nabla\psi_\theta(x_b)\)
  \State Compute the declared dual or potential-based training objective, including any required convex-conjugate or surrogate terms
  \State Add any declared regularization enforcing smoothness, stability, or numerical tractability
  \State Backpropagate the total loss and update \(\theta\)
  \State Re-enforce convexity according to the declared ICNN policy if it is not built directly into the parameterization
\EndWhile
\State On held-out validation clouds, compute mapped particles \(T_\theta(x_i)\)
\State Evaluate endpoint pushforward mismatch, transport-cost proxy, and downstream scalar sensitivity on the held-out clouds
\State Return \(\psi_\theta\), the induced map \(T_\theta\), and the declared diagnostics
\end{algorithmic}
\end{algorithm}

\paragraph{Approximation enters here.}
The learned object is the transport map itself.  Errors in the learned potential
or its gradient directly change the transported output cloud.

\section{Monge-gap route: flexible direct-map implementation}
\label{sec:bf-neural-ot-direct-monge-gap}

A more flexible direct-map family regularizes a candidate map by how far it
behaves from an OT map rather than forcing a particular architecture from the
start.  For cost \(c\) and reference measure \(\rho\), the Monge gap is
\begin{equation}
\label{eq:bf-neural-ot-direct-monge-gap}
  \mathcal M_\rho^c(T)
  =
  \int c(x,T(x))\,d\rho(x)-W_c(\rho,T_\#\rho).
\end{equation}
As a mathematical definition, this is exact.  But in a practical training loop,
the second term is typically not available in closed form because the pushforward
measure \(T_\#\rho\) itself depends on the current learned map.  So any finite-
sample training objective that uses an estimated OT reference term should be read
as a surrogate or approximate Monge-gap objective unless the estimator and its
properties are stated explicitly.  Operationally, the chapter’s point is: fit a
useful map, but penalize it when it behaves unlike an OT map for the chosen
cost.

On the running cell, the good and bad endpoint maps above already show the idea.
Both endpoint laws can match, but the bad map moves mass much less efficiently.
The Monge-gap route is attractive precisely because it can discourage such maps
without forcing the full ICNN/Brenier structural commitment.

\begin{algorithm}[ht]
\caption{\texttt{train\_monge\_gap\_regularized\_direct\_map}}
\begin{algorithmic}[1]
\Require source-cloud sampler for \(\mu\), target-cloud sampler for \(\nu\), reference-measure sampler for \(\rho\), parametric map family \(T_\theta\), cost function \(c\), endpoint-fit objective, Monge-gap weight \(\lambda_{\mathrm{MG}}\), optimizer, stopping rule, and held-out validation suite
\Ensure trained direct map \(T_\theta\), Monge-gap diagnostic history, endpoint-fit diagnostics, and downstream scalar sensitivity report
\State Initialize map parameters \(\theta\) and optimizer state
\While{the stopping rule has not been met}
  \State Draw minibatches from the source, target, and reference distributions
  \State Compute mapped source particles \(z_i=T_\theta(x_i)\)
  \State Evaluate the declared endpoint-fit loss between the mapped source batch and the target batch
  \State Compute the empirical transport-cost term on the reference samples from \(\rho\)
  \State Estimate the OT reference term needed for the empirical Monge-gap penalty; this estimate defines the practical training surrogate unless a stronger equivalence result is supplied
  \State Form the total loss as endpoint-fit term plus \(\lambda_{\mathrm{MG}}\) times the empirical Monge-gap penalty, together with any declared regularization
  \State Backpropagate the total loss and update \(\theta\)
\EndWhile
\State On held-out validation clouds, evaluate endpoint mismatch, Monge-gap magnitude, transport-cost proxy, and downstream scalar drift
\State Return the trained map \(T_\theta\) and the declared diagnostics
\end{algorithmic}
\end{algorithm}

\paragraph{Reader checkpoint.}
The Monge-gap route still learns a direct map, but it relaxes the strong
ICNN-only structural commitment.  The gain is flexibility.  The cost is that
the learned map is now judged relative to a regularized OT-like behavior rather
than by a retained teacher solver.

\section{Structured variants and why this is not the default first route}
\label{sec:bf-neural-ot-direct-structured}
\label{sec:bf-neural-ot-direct-why-not-default}

The same direct-map logic can be combined with additional structural priors, such
as sparse displacement or factorized transport structure.  These are worth
keeping in view because they may be useful when the filtering update itself has a
scientifically justified sparse or low-active-subspace geometry.  But they
inherit the same chapter-wide burden: the learned endpoint map itself must be
validated, rather than being interpreted through a retained corrective teacher.

That is also why direct-map learning is not the monograph’s default immediate
route.  It is mathematically serious and potentially powerful, but it learns the
final map object itself rather than only accelerating a declared teacher solve,
and its strongest theorem-level support depends on more specific geometric
assumptions than the retained-teacher route requires.

\section{Verification and non-claims}
\label{sec:bf-neural-ot-direct-verification}

A minimal implementation-facing reading of this chapter should distinguish:
\begin{itemize}
  \item what map object is being fit,
  \item what assumptions justify that map family,
  \item what endpoint, pushforward, transport-cost, or Monge-gap residuals are
  measured,
  \item what downstream scalar drift is checked once the learned map is embedded
  in filtering or HMC-adjacent code.
\end{itemize}

The practical contract differs from the retained-teacher chapter because no
corrective Sinkhorn phase is automatically present to return to a declared
teacher object.  So endpoint fit, pushforward mismatch, transport-cost proxy,
Monge-gap proxy, and downstream scalar sensitivity are part of the primary
evidence, not side diagnostics.

This chapter does \emph{not} claim:
\begin{itemize}
  \item that direct learned maps are interchangeable with retained-teacher EOT
  acceleration,
  \item that map-level fit automatically preserves the same filtering scalar,
  \item that broader direct-map families are already validated for BayesFilter.
\end{itemize}
Its role is to slow down the direct-map family itself so the reader can see that
it is solving a different problem from the previous chapter.

\FloatBarrier
